%% kapa: exam=1 round=34 year=2023 period=1 form=A subject=civil_law subject_ko=민법 seq=15 number=15
%% source: Q-Net 감정평가사 기출문제 (cst00309) artlSeq=5211789 file=raw/1008/5211789/2023년도 제34회 감정평가사 1차시험 1교시.pdf
%% conversion: kordoc 4.24.51 -> tools/gen1.py (auto); review=proofread-sonnet; flags=none
\begin{question}{15}
복대리에 관한 설명으로 옳은 것은? (다툼이 있으면 판례에 따름)
\begin{choices}
  \item 복대리인은 대리인의 대리인이다.
  \item 복대리인은 본인에 대해 어떠한 권리·의무도 부담하지 않는다.
  \item 복대리인이 선임되면 복대리인의 대리권 범위 내에서 대리인의 대리권은 잠정적으로 소멸한다.
  \item 대리인이 복대리인을 선임한 후 사망하더라도 특별한 사정이 없는 한 그 복대리권은 소멸하지 않는다.
  \item 복임권 없는 대리인에 의해 선임된 복대리인의 대리행위에 대해서도 권한을 넘은 표현대리에 관한 규정이 적용될 수 있다.
\end{choices}
\end{question}
\answer{5}
